Este informe presenta los resultados de la implementación y análisis experimental de distintos algoritmos de ordenamiento (MergeSort, QuickSort, y Sort) y de multiplicación de matrices (Naive y Strassen). Se realiza un análisis comparativo entre los algoritmos, considerando para ello el tiempo empleado para ejecutar un algoritmo, y la memoria ocupada, esto para distintos tipos y tamaños de entrada. Como resultado experimental, se genera gráficos para cada algoritmo, conforme a cada tipo de entrada. Según los gráficos resultantes, entre los algoritmos de ordenamiento, "Sort", de la librería estándar de C++ es el más eficiente porque cuenta con los menores tiempos de ejecución. Para los algoritmos de multiplicación de matrices, "Naive" tiene ventaja, ya que, para todos los casos de prueba, fue más rápido. Se concluye que, para escoger el mejor algoritmo, hay que considerar factores como la complejidad asintótica, pero también, hay que tener en cuenta las características de hardware, y los patrones de los casos de prueba de entrada. Además, se sugiere que para mejorar el rendimiento se debe estudiar las implementaciones híbridas de algoritmos.

